% \documentclass[aspectratio=169,notes]{beamer}
\documentclass[aspectratio=169]{beamer}
\usetheme[faculty=phil]{fibeamer}
\usepackage{polyglossia}
\setmainlanguage{english} %% main locale instead of `english`, you
%% can typeset the presentation in either Czech or Slovak,
%% respectively.
\setotherlanguages{russian} %% The additional keys allow
%%
%%   \begin{otherlanguage}{czech}   ... \end{otherlanguage}
%%   \begin{otherlanguage}{slovak}  ... \end{otherlanguage}
%%
%% These macros specify information about the presentation
\title[Artificial Intelligence Software (AIS)]{Artificial Intelligence Software, HW 7} %% that will be typeset on the
\subtitle{Model Deployment with Triton
\\ \  \\ \  } %% title page.
\author{Oleg Bulichev}
%% These additional packages are used within the document:
\usepackage{ragged2e}  % `\justifying` text
\usepackage{booktabs}  % Tables
\usepackage{tabularx}
\usepackage{tikz}      % Diagrams
\usetikzlibrary{calc, shapes, backgrounds}
\usetikzlibrary{decorations.pathreplacing,calligraphy,calc,graphs}
\usepackage{amsmath, amssymb}
\usepackage{url}       % `\url`s
\usepackage{listings}  % Code listings
% \usepackage{subfigure}
\usepackage{floatrow}
\usepackage{subcaption}
\usepackage{mathtools}
\usepackage{todonotes}
\usepackage{fontspec}
\usepackage{multicol}
\usepackage{pdfpages}
\usepackage{wrapfig}
\usepackage{qrcode}
\usepackage{animate}
\usepackage{booktabs}
\usepackage{multirow}
% \usepackage{graphicx}
\usepackage{colortbl}

\graphicspath{{resources/}}
\frenchspacing

\setbeamertemplate{caption}[numbered]
\usetikzlibrary{graphs}

% \usepackage[backend=biber,style=ieee,autocite=footnote]{biblatex}
% \addbibresource{biblio.bib}
% \DefineBibliographyStrings{english}{%
%   bibliography = {References},}

\newcommand{\oleg}[2][] {\todo[color=red, #1] {OLEG:\\ #2}}
\newcommand{\fbckg}[1]{\usebackgroundtemplate{\includegraphics[width=\paperwidth]{#1}}}%frame background

\usepackage[framemethod=TikZ]{mdframed}
\newcommand{\dbox}[1]{
\begin{mdframed}[roundcorner=3pt, backgroundcolor=yellow, linewidth=0]
\vspace{1mm}
{#1}
\vspace{1mm}
\end{mdframed}
}

\begin{document}
\setlength{\abovedisplayskip}{0pt}
\setlength{\belowdisplayskip}{0pt}
\setlength{\abovedisplayshortskip}{0pt}
\setlength{\belowdisplayshortskip}{0pt}

\fbckg{fibeamer/figs/title_page.png}
\frame[c]{\setcounter{framenumber}{0}
    \usebeamerfont{title}%
    \usebeamercolor[fg]{title}%
    \begin{minipage}[b][6.5\baselineskip][b]{\textwidth}%
        \textcolor{black}{\raggedright\inserttitle}
    \end{minipage}
    % \vskip-1.5\baselineskip

    \usebeamerfont{subtitle}%
    \usebeamercolor[fg]{framesubtitle}%
    \begin{minipage}[b][3\baselineskip][b]{\textwidth}
        \raggedright%
        \insertsubtitle%
    \end{minipage}
    \vskip.25\baselineskip
}
%   \frame[c]{\maketitle}

\fbckg{fibeamer/figs/common.png}

\note{\scriptsize \begin{itemize}
        \item \
    \end{itemize}}

\begin{frame}[t]{Task 1 --- Model Preparation \& Triton Deployment}
\vspace{-0.3cm}
\textbf{Short description}: Deploy a YOLO instance segmentation model to production using the full NVIDIA stack, following the tutorial below.

\textbf{Artifacts}: Repository with code and PDF report with screenshots for each step.

\textbf{Requirements}
\footnotesize
Tutorial \& repo: \url{https://youtu.be/ljqyuDxd_H0} \quad \url{https://github.com/Koldim2001/Triton_example}\\[2pt]
Use \texttt{yolov8n-seg} (instance segmentation) instead of the ResNet-152 shown in the tutorial.
\begin{enumerate}
    \item \textbf{Export to ONNX}: export the model from PyTorch with a dummy input and dynamic axes so it accepts variable batch sizes.
    \item \textbf{Triton file structure}: organize \texttt{models/<name>/<version>/model.onnx} and write \texttt{config.pbtxt} (platform, input/output tensor names, types, shapes, \texttt{max\_batch\_size}).
    \item \textbf{Launch Triton in Docker} with GPU passthrough. Expose ports \texttt{8000} (HTTP), \texttt{8001} (gRPC), \texttt{8002} (metrics).
    \item \textbf{FastAPI async service}: implement preprocessing identical to training pipeline, call Triton via \texttt{tritonclient} over gRPC asynchronously, return segmentation results (masks + class labels).
\end{enumerate}

\end{frame}

\begin{frame}[t]{Task 1 --- TensorRT, Monitoring \& Advanced Triton}
\footnotesize
\begin{enumerate}
    \setcounter{enumi}{4}
    \item \textbf{TensorRT optimization}: convert the ONNX model to \texttt{.plan} using \texttt{trtexec} inside the NVIDIA container. Enable FP16 quantization for maximum throughput on the target GPU.
    \item \textbf{Monitoring}: configure Prometheus to scrape metrics from port \texttt{8002}; import a Grafana dashboard showing RPS, latency breakdown (preprocessing / inference / postprocessing), GPU utilization, and average batch size.
    \item \textbf{Advanced Triton features}:
    \begin{itemize}
        \item Enable \textbf{dynamic batching} to merge concurrent single-image requests into one GPU batch.
        \item Configure \textbf{instance group} to run multiple model copies in parallel.
        \item Run a load test (e.g.\ Locust) and show Grafana metrics before and after enabling these features.
    \end{itemize}
\end{enumerate}

\end{frame}


\fbckg{fibeamer/figs/last_page.png}
\frame[plain]{}

\end{document}